\chapter{مقدمه}
\label{ch:introduction}

\section{مسئله اعتماد به شبکه‌های عصبی}
شبکه‌های عصبی عمیق در دسته‌بندی تصویر، پردازش زبان و تحلیل داده‌های پیچیده به ابزارهایی پرکاربرد تبدیل شده‌اند. بااین‌حال، درستی زیاد روی نمونه‌های آزمون معمول لزوماً به معنای پایداری تصمیم در همسایگی همان نمونه‌ها نیست. پژوهش‌های آغازین نشان دادند می‌توان با تغییری کوچک و هدفمند در ورودی، خروجی یک شبکه را تغییر داد؛ تغییری که گاه برای مشاهده‌گر انسانی به‌آسانی تشخیص‌پذیر نیست \cite{szegedy2014intriguing,goodfellow2015explaining}. چنین ورودی‌ای نمونه هم‌آورد نامیده می‌شود. در این نوشتار، برابر «هم‌آورد» برای \lr{adversarial} و «آشفتگی» برای \lr{perturbation} به کار می‌رود. عنوان پایان‌نامه مطابق پیشنهاد اولیه حفظ شده است و «تخاصمی» در آن به همین مفهوم اشاره دارد.

آسیب‌پذیری هم‌آورد یک پرسش بنیادی درباره آنچه مدل از داده آموخته است مطرح می‌کند: آیا تصمیم مدل در یک ناحیه معنادار پیرامون نمونه پایدار است یا تنها در خود نقاط آزمون درست به نظر می‌رسد؟ این پرسش به‌ویژه هنگامی مهم می‌شود که خروجی مدل بخشی از یک سامانه تصمیم‌گیری باشد. البته وجود یک دفاع تجربی در برابر آشفتگی‌های محدود پیکسلی به‌تنهایی ایمنی یک سامانه پزشکی یا خودمختار را اثبات نمی‌کند؛ ارزش این چارچوب، فراهم‌کردن مسئله‌ای دقیق و سنجش‌پذیر برای مطالعه بخشی از آسیب‌پذیری است.

در میان دفاع‌های پیشنهادی، آموزش هم‌آورد جایگاه ویژه‌ای دارد. در این روش، شبکه تنها با ورودی پاک آموزش نمی‌بیند؛ نمونه‌هایی که برای افزایش هزینه مدل ساخته شده‌اند نیز در فرایند یادگیری حضور دارند. هدف آن است که مدل در دشوارترین نقاط یک مجموعه آشفتگی ازپیش‌تعریف‌شده نیز عملکرد قابل‌قبولی داشته باشد \cite{madry2018towards}. مشکل عملی آن است که یافتن این نقاط دشوار معمولاً به چند بار محاسبه شیب نسبت به ورودی نیاز دارد و هزینه آموزش را افزایش می‌دهد.

\section{از آموزش چندگامی تا آموزش سریع}
صورت متداول آموزش هم‌آورد یک مسئله کمینه‌سازی و بیشینه‌سازی تودرتو است. مسئله درونی، ورودی نامطلوب را برای پارامترهای فعلی می‌جوید و مسئله بیرونی پارامترها را با توجه به آن بهبود می‌دهد. در روش چندگامی، حمله‌ای مانند \lr{PGD} چند بار در جهت افزایش هزینه حرکت می‌کند و پس از هر حرکت به مجموعه مجاز تصویر می‌شود. این تکرار، تقریب قوی‌تری از مسئله درونی فراهم می‌کند، اما زمان و منابع بیشتری می‌طلبد.

آموزش هم‌آورد سریع با جایگزین‌کردن حمله چندگامی با حمله تک‌گامی، این هزینه را کاهش می‌دهد \cite{wong2020fast}. جذابیت این ایده روشن است: اگر یک گام مبتنی بر شیب بتواند نمونه‌ای به‌اندازه کافی دشوار بسازد، بخش عمده هزینه حمله حذف می‌شود. بااین‌حال، کیفیت این تقریب به هندسه تابع هزینه وابسته است. یک گام ممکن است در ناحیه‌ای تقریباً خطی مناسب باشد، اما پس از تغییر شکل تابع هزینه از یک ناحیه آسیب‌پذیر عبور کند و به نقطه‌ای با هزینه کمتر برسد.

بنابراین، «سریع» بودن تنها با شمردن گام‌های حمله تعریف نمی‌شود. روشی که یک اجرای آن ارزان است ولی برای یافتن تنظیم پایدار به چندین اجرای ناموفق نیاز دارد، ممکن است در عمل هزینه زیادی داشته باشد. همچنین، کاهش زمان همراه با سقوط درستی پاک یا شکست روی یک دادگان دیگر، پاسخ کاملی به مسئله آموزش قابل‌اتکا نیست. در این پژوهش، هزینه هر اجرا، پایداری، درستی پاک و درستی هم‌آورد هم‌زمان بررسی می‌شوند.

\section{بیش‌برازش فاجعه‌بار و شکاف ارزیابی}
بیش‌برازش فاجعه‌بار به حالتی گفته می‌شود که در طول آموزش سریع، درستی مدل در برابر حملات چندگامی در زمانی کوتاه به‌شدت افت می‌کند، درحالی‌که درستی در برابر حمله تک‌گامی آموزش ممکن است همچنان زیاد یا حتی رو به افزایش باشد \cite{wong2020fast,andriushchenko2020understanding}. این ناهمخوانی نشان می‌دهد مدل در برابر روش خاص تولید نمونه آموزشی سازگار شده است، اما این سازگاری به تمام ناحیه تهدید تعمیم پیدا نکرده است.

ویژگی خطرناک پدیده آن است که شاخص‌های ظاهراً مطلوب آموزش می‌توانند هم‌زمان با ازبین‌رفتن ستبری دیده شوند. کاهش هزینه روی نمونه‌های تولیدشده با \lr{FGSM} لزوماً به معنای کاهش بیشینه هزینه در کره آشفتگی نیست. اگر خود مولد نمونه هم‌آورد شکست خورده باشد، بهینه‌سازی بیرونی روی مسئله‌ای آسان‌تر از مسئله واقعی ادامه می‌یابد. ارزیابی با همان حمله آموزش، این شکست را پنهان می‌کند.

این وضعیت دو نیاز مرتبط ایجاد می‌کند. نخست، باید نشانه‌ای از نامعتبرشدن تقریب تک‌گامی را پیش از فروپاشی آشکار درستی چندگامی مشاهده کرد. دوم، باید از اطلاعات همان نشانه برای اصلاح ساخت نمونه آموزشی استفاده کرد. پایش دائمی یک حمله قوی راهی مستقیم برای مشاهده افت است، اما ممکن است با هدف کم‌هزینه‌بودن آموزش سریع ناسازگار شود. مسئله اصلی پایان‌نامه، دستیابی به اطلاعات هشداردهنده از محاسباتی است که از قبل در آموزش انجام می‌شوند.

\section{پرسش‌ها و اهداف پژوهش}
پرسش نخست این است که محدودبودن اندازه و جهت آشفتگی آموزش چگونه با تغییر هندسه تابع هزینه و بیش‌برازش فاجعه‌بار مرتبط می‌شود. برای پاسخ، باید رفتار مدل را نه‌فقط در اندازه آشفتگی آموزش، بلکه در نقاط میانی همان جهت نیز بررسی کرد. اگر مدل تنها در چند فاصله خاص عملکرد مطلوبی داشته باشد، درستی بالای حمله آموزش می‌تواند تصویری ناقص از ستبری ارائه دهد.

پرسش دوم به پیش‌بینی مربوط است: آیا مقایسه شیب ورودی پیش و پس از حمله، افت کیفیت تقریب تک‌گامی را پیش از جدایی آشکار درستی تک‌گامی و چندگامی نشان می‌دهد؟ پاسخ کاربردی باید هم از نظر هندسی تفسیرپذیر باشد و هم برای محاسبه به گذرهای متعدد اضافی نیاز نداشته باشد. ارتباط با خمیدگی مهم است، اما محاسبه کامل هسین در ابعاد تصاویر گزینه مناسبی برای روش سریع نیست.

پرسش سوم، ساخت یک سازوکار اصلاحی است. اگر از اختلاف دو شیب بتوان اطلاعاتی درباره خمیدگی جهتی به دست آورد، آیا می‌توان اندازه گام حمله را با آن تنظیم کرد؟ چنین سازوکاری باید نوسان دسته‌های آموزشی را نیز تحمل کند و از وابستگی شدید به تنظیم جداگانه برای هر دادگان یا معماری بکاهد. ارزیابی تجربی باید نشان دهد این هدف در چه تنظیم‌هایی برآورده شده و در چه مواردی محدودیت باقی مانده است.

\section{دستاوردها و نسبت آن‌ها با مقاله مشترک}
هسته علمی پایان‌نامه بر پژوهش مشترک با عنوان
\lr{SORA: Free Second-Order Attacks in Fast Adversarial Training}
استوار است که در \lr{ICML 2026} منتشر شده است \cite{sora2026}. نویسندگان مقاله، به ترتیب، \lr{Mazdak Teymourian}، \lr{Ramtin Moslemi}، فرزان رحمانی و دکتر محمدحسین رهبان هستند. مطالب مربوط به آزمایش‌ها و روش \sora{} در این نوشتار به این کار مشترک ارجاع داده می‌شوند؛ ترتیب نویسندگان، به‌تنهایی مبنایی برای نسبت‌دادن بخش خاصی از پژوهش به یک فرد نیست.

سه محور اصلی پژوهش عبارت‌اند از:
\begin{enumerate}
\item بررسی بیش‌برازش اپسیلون و نشان‌دادن اهمیت پوشش اندازه‌ها و جهت‌های متنوع آشفتگی در آموزش سریع؛
\item تعریف شاخص هم‌راستایی آشفتگی یا \pertalign{} با استفاده از دو شیب موجود در فرایند آموزش و تحلیل ارتباط محلی آن با هسین؛
\item طراحی روش \sora{} برای انتخاب تطبیقی توزیع اندازه گام و ارزیابی آن روی دادگان و معماری‌های گوناگون.
\end{enumerate}

گزارش سمینار پیشین، مسئله و کارهای مرتبط را مرور کرده و ایده هم‌راستایی آشفتگی و تنظیم گام را در مرحله اولیه مطرح کرده بود. پیشنهاد پایان‌نامه نیز بر پیش‌بینی کم‌هزینه بیش‌برازش تمرکز داشت. مقاله نهایی، این مسیر را با تحلیل و ارزیابی گسترده‌تر تکمیل کرده است. بنابراین، ایده‌های آینده در متن سمینار در این پایان‌نامه صرفاً زمانی به‌عنوان نتیجه ارائه می‌شوند که مقاله برای آن‌ها شواهد مشخص داشته باشد.

\section{دامنه و حدود ادعاها}
دامنه اصلی مطالعه، دسته‌بندی تصویر و آشفتگی‌های محدود با نرم $\ell_\infty$ است. نتایج اصلی بر ارزیابی‌های تجربی با حملات مشخص متکی هستند و گواهی ریاضی ستبری برای تمام ورودی‌های مجاز محسوب نمی‌شوند. همچنین، تفسیر مرتبه دوم بر تقریب محلی و فرض‌های همواری متکی است؛ در یک شبکه واقعی با تابع‌های فعال‌سازی تکه‌ای و برش پیکسل، این تفسیر باید با احتیاط به کار رود.

تمام اعداد گزارش‌شده در فصل آزمایش‌ها از نسخه مقاله و پیوست ارائه‌شده استخراج شده‌اند. در آماده‌سازی این نوشتار، آموزش شبکه یا آزمایش آماری تازه‌ای اجرا نشده است. هرجا متن اصلی و پیوست، عدد یا تنظیم متفاوتی ارائه می‌کنند، محل گزارش مشخص می‌شود. این تفکیک برای جلوگیری از ترکیب نادرست آزمایش‌ها و برای بازسازی دقیق‌تر نتایج ضروری است.

همچنین، مشاهده هشدار زودتر در یک نمودار با ساخت یک پیش‌بین کالیبره‌شده یکسان نیست. تعیین آستانه عمومی، نرخ هشدار کاذب و توزیع زمان پیش‌آگاهی به ارزیابی جداگانه نیاز دارد. پایان‌نامه، شواهد موجود درباره هشدار را تفسیر می‌کند و این ارزیابی‌های تکمیلی را به‌عنوان مسیر ادامه پژوهش مشخص می‌سازد.

\section{ساختار پایان‌نامه}
فصل دوم مفاهیم دسته‌بندی، مدل تهدید، حملات و آموزش هم‌آورد را تعریف می‌کند. فصل سوم کارهای مرتبط را بر پایه سازوکار مقابله با بیش‌برازش مرور می‌کند. فصل چهارم به بیش‌برازش اپسیلون و شاخص \pertalign{} اختصاص دارد. در فصل پنجم، استخراج اندازه گام مرتبه دوم و الگوریتم عملی \sora{} ارائه می‌شود. فصل ششم نتایج، تنظیم آزمایش‌ها و محدودیت‌های مقایسه را بررسی می‌کند. فصل هفتم به جمع‌بندی و پیشنهادهای آینده می‌پردازد. پیوست، نمادها و اطلاعات لازم برای دنبال‌کردن منبع و بازتولید پژوهش را گردآوری می‌کند و واژه‌نامه، معادل‌های فارسی و انگلیسی اصطلاحات را در دسترس قرار می‌دهد.
